\section{Información como recurso, y tres precisiones necesarias}
\label{sec:informacion2}

Las secciones anteriores tratan la información como restricción —quién habla con
quién— pero no como \emph{recurso} que la coalición pueda adquirir. Esta sección
cierra esa asimetría y, de paso, precisa tres puntos donde los enunciados
previos admiten una lectura más fuerte de la que sostienen.

\subsection{Reducir incertidumbre amplía lo certificable}

Los parámetros de la carga —masa $m_\ell$, inercia $I_\ell$, coeficiente de
fricción— rara vez se conocen con exactitud. Sea $\Theta$ el conjunto de
parámetros compatibles con lo observado y
\begin{equation}
  \Kcal(d,X;\Theta)\defeq\bigcap_{\theta\in\Theta}\Kcal(d,X;\theta)
  \label{eq:factibilidad-robusta}
\end{equation}
la factibilidad robusta de un compromiso $d$: las acciones válidas para
\emph{todo} parámetro compatible.

\begin{proposicion}[Monotonía de la información]
\label{pr:monotonia-info}
A estado y problema congelados, $\Theta_2\subseteq\Theta_1$ implica
$\Kcal(d,X;\Theta_1)\subseteq\Kcal(d,X;\Theta_2)$.
\end{proposicion}

\begin{proof}
Una acción válida para todo $\theta\in\Theta_1$ lo es en particular para todo
$\theta$ de cualquier subconjunto. La inclusión de conjuntos invierte el sentido
al intersecar sobre menos índices.
\end{proof}

El enunciado es elemental y su consecuencia no lo es: \textbf{reducir incertidumbre válida amplía el conjunto de
continuaciones certificables sin añadir un solo AMR}. Con velocidad $1$~m/s y fuerza de frenado $60$~N, una masa desconocida en
$[50,150]$~kg obliga a reservar $1{,}25$~m de distancia de parada; si la
observación estrecha el intervalo a $[70,90]$~kg, bastan $0{,}75$~m. La
información no crea capacidad física: elimina conservadurismo.

\begin{corolario}[Medir es una acción del juego]
\label{co:medir}
La parte discreta $d_i$ de \cref{def:estrategia} admite una acción adicional
—ejecutar una maniobra informativa— cuyo pago es la ampliación de
$\Kcal(d,X;\Theta)$ que produce, y cuyo coste es el tiempo y el riesgo de la
maniobra. La decisión «reclutar otro AMR» y la decisión «demostrar que los
actuales bastan» compiten en el mismo potencial.
\end{corolario}

\begin{observacion}[Seguir bien no es haber identificado]
\label{obs:identificacion}
Que la carga siga la trayectoria con error pequeño no implica que sus parámetros
sean identificables. Con el regresor planar
$\bigl[F_x,F_y,T_z\bigr]^{\!\top}=Y_\ell\,[m_\ell,I_\ell]^{\!\top}+\epsilon$ y
maniobras puramente rectilíneas, la columna asociada a $I_\ell$ es nula y la
inercia angular no se identifica por bien que se siga la recta. Análogamente,
$|f_t|\le\mu N$ durante adherencia da una \emph{cota inferior} de $\mu$, no su
valor. Una maniobra informativa debe además ser físicamente segura: su coste y
su riesgo entran en la misión, y no puede añadirse un término de exploración
suponiendo que todo movimiento informativo está permitido.
\end{observacion}

\subsection{Precisión 1: el orden de los cuantificadores}

\Cref{th:robusto} afirma $S_W\Wdem+\delta\mathbb B\subseteq\mathcal W(\Ccal)$, es
decir: \emph{para toda} perturbación admisible \emph{existe} un reparto que la
realiza. Ese es el orden $\forall\theta\,\exists u$, correcto cuando la
perturbación se \emph{mide} y el programa se resuelve de nuevo con ella.

No es el orden que hace falta cuando el reparto debe fijarse antes de conocer la
perturbación. Entonces se requiere $\exists u\,\forall\theta$, estrictamente más
fuerte y en general no implicado por el anterior.

\begin{observacion}[Los dos órdenes no son equivalentes]
\label{obs:cuantificadores}
Con $u\in[-2,2]$ y $\theta\in\{-1,+1\}$, la restricción $\theta u\ge1$ es
factible para cada $\theta$ por separado —$u=-2$ sirve para $\theta=-1$ y $u=2$
para $\theta=+1$— y ningún $u$ la satisface para ambos. Certificar una coalición
comprobando cada escenario con un control distinto produce un certificado falso
si el escenario no será conocido en el momento de actuar.
\end{observacion}

\subsection{Precisión 2: un grafo dirigido acuerda, pero acuerda mal}

\Cref{th:mapa} supone grafo no dirigido, y la hipótesis no es cosmética. Bajo
$\dot x=-Lx$ con digrafo fuertemente conexo, el estado converge a
$\mathbf 1\pi^{\!\top}x(0)$ con $\pi^{\!\top}L=0$; la media aritmética se
conserva solo si el grafo es balanceado en peso.

\begin{observacion}[Acuerdo sobre el valor equivocado]
\label{obs:digrafo}
Con
$L=\bigl[\begin{smallmatrix}1&-1&0\\0&1&-1\\-2&0&2\end{smallmatrix}\bigr]$ se
obtiene $\pi=(0{,}4,\,0{,}4,\,0{,}2)$; partiendo de $x(0)=(0,0,9)$ los tres
agentes convergen a $1{,}8$, mientras la media es $3$. Todos coinciden y todos
se equivocan. Si ese valor alimenta la minimización de
$\Phi(z)=\tfrac12\sum_i(z-a_i)^2$, el exceso de coste es $2{,}16$.
La consecuencia para \cref{th:mapa} y para el acuerdo de precios de
\cref{th:precio-seguridad} es que \emph{no basta comprobar que las copias convergen}: hay que demostrar que el agregado al que convergen es el que el
objetivo social requiere. El árbol generador de \cref{th:topologia-garantizada}
debe entenderse, por tanto, como subgrafo no dirigido.
\end{observacion}

\subsection{Precisión 3: el mercado no es un descenso, es un punto de silla}

\Cref{th:muestreada} presenta la actualización de precios como descenso de
gradiente sobre la función dual $F$, y como tal converge para
$0<\gamma<2/L$. Esa lectura es válida cuando $\nabla F$ se evalúa de forma
exacta, es decir, resolviendo el problema primal interior hasta optimalidad en
cada paso.

Una implementación distribuida rara vez hace eso: actualiza primal y dual
simultáneamente. El operador conjunto
\begin{equation}
  F(u,\lambda)=
  \begin{pmatrix}\nabla\Phi(u)+A^{\!\top}\lambda\\ b-Au\end{pmatrix}
  \label{eq:operador-silla}
\end{equation}
tiene parte antisimétrica y \emph{no} es el gradiente de ninguna función en el
par de variables.

\begin{proposicion}[El paso de Euler simultáneo diverge]
\label{pr:euler-diverge}
Para la silla $\mathcal L(x,\lambda)=x\lambda$, el paso explícito
$z^{k+1}=z^k-hF(z^k)$ satisface $\|z^{k+1}\|^2=(1+h^2)\|z^k\|^2$, luego diverge
para todo $h>0$. El paso extragradiente
$y^k=z^k-hF(z^k)$, $z^{k+1}=z^k-hF(y^k)$ satisface
$\|z^{k+1}\|^2=(1-h^2+h^4)\|z^k\|^2$, que decrece para $0<h<1$.
\end{proposicion}

\begin{proof}
Con $F(x,\lambda)=(\lambda,-x)$, el paso explícito da
$z^{+}=\bigl((x-h\lambda),(\lambda+hx)\bigr)$ y desarrollando el cuadrado los
términos cruzados se cancelan, quedando $(1+h^2)\|z\|^2$. Para el
extragradiente, $F(y)=(\lambda+hx,\,-x+h\lambda)$ y
$z^{+}=\bigl((1-h^2)x-h\lambda,\ hx+(1-h^2)\lambda\bigr)$, cuyo cuadrado es
$\bigl[(1-h^2)^2+h^2\bigr]\|z\|^2=(1-h^2+h^4)\|z\|^2$.
\end{proof}

Numéricamente, con $h=0{,}4$ y norma inicial unidad, tras sesenta iteraciones el
paso explícito alcanza $85{,}85$ y el extragradiente $1{,}3\times10^{-2}$.

\begin{observacion}[Consecuencia de diseño]
\label{obs:extragradiente}
Antes de atribuir una oscilación de precios a la física del contacto o a la
dinámica poblacional conviene auditar el esquema de integración del lazo
primal--dual. El extragradiente no es una aportación de este trabajo: es una
familia establecida para desigualdades variacionales. Lo que aquí se afirma es
que \cref{th:muestreada} cubre el caso de gradiente dual exacto y \emph{no} el
de actualización simultánea, y que en ese segundo caso hace falta un esquema con
paso predictor--corrector.
\end{observacion}
